\documentclass[11pt, letterpaper]{article}

% --- PAQUETES ---
\usepackage[utf8]{inputenc}
\usepackage[T1]{fontenc}
\usepackage[spanish, es-noshorthands]{babel}
\usepackage[top=2cm, bottom=2.2cm, left=2.2cm, right=2.2cm]{geometry}
\usepackage{graphicx}
\usepackage{amsmath}
\usepackage{amssymb}
\usepackage{enumitem}
\usepackage[table]{xcolor}
\usepackage{tikz}
\usetikzlibrary{arrows.meta}
\usepackage{fancyhdr}
\usepackage{eso-pic}
\usepackage[most]{tcolorbox}
\usepackage{helvet}
\renewcommand{\familydefault}{\sfdefault}

% --- COLORES INSTITUCIONALES (LICEO V.A.L.) ---
\definecolor{valblue}{RGB}{0,112,192}
\definecolor{valdark}{RGB}{30,30,30}
\definecolor{valgray}{RGB}{245,247,250}

% Fondo institucional obligatorio del Liceo V.A.L.
\AddToShipoutPictureBG{
	\AtPageLowerLeft{
		\includegraphics[width=\paperwidth,height=\paperheight,keepaspectratio=false]{../../FondoVAL.pdf}
	}
}

\begin{document}
	
	\begin{center}
		{\Large \textbf{\color{valblue} LICEO V.A.L. (Vida, Amor, Luz)}}\\[0.1cm]
		{\large \textbf{DEPARTAMENTO DE CIENCIAS BÁSICAS}}\\[0.1cm]
		{\large \textbf{EXAMEN DE FÍSICA: NATURALEZA DE LA
				LUZ. REFLEXIÓN}}\\[0.1cm]
		{\small \textbf{Física 11$^{\circ}$} \quad | \quad \textbf{Código: 11F -- 05} \quad | \quad \textbf{CON CALCULADORA}}
	\end{center}
	
	\vspace{0.2cm}
	\noindent \textbf{NOMBRE:} \underline{\hspace{8.5cm}} \hfill \textbf{FECHA:} \underline{\hspace{3cm}}
	\vspace{0.3cm}
	
	
	
	\begin{enumerate}[leftmargin=*, itemsep=0.8cm]
		
		% PUNTO 1
		\item \textbf{(25\%)}
		Se tiene un espejo cóncavo con una distancia focal de $15\text{ cm}$.
		\begin{enumerate}[label=\alph*.), itemsep=0.3cm]
			\item ¿A qué distancia se debe colocar un objeto de $5\text{ cm}$ de altura para que su imagen se forme a $30\text{ cm}$ del espejo?
			\item ¿Cuál es la altura de la imagen?
		\end{enumerate}
		
	
			\small			
				\[\frac{1}{f} = \frac{1}{d_o} + \frac{1}{d_i} \qquad|\qquad A = \frac{h_i}{h_o} = -\frac{d_i}{d_o} \]
		
		
		% PUNTO 2
		\item \textbf{(25\%)} A $8\text{ cm}$ de un espejo convexo de $10\text{ cm}$ de radio, se coloca un objeto de $4\text{ cm}$ de altura.
		\begin{enumerate}[label=\alph*.), itemsep=0.3cm]
			\item ¿A qué distancia se forma la imagen?
			\item ¿Cuál es el tamaño de la imagen?
		\end{enumerate}
		
	
			\small
			\[f = \frac{R}{2}  \qquad|\qquad  \frac{1}{f} = \frac{1}{d_i} - \frac{1}{d_o} \qquad|\qquad A = \frac{h_i}{h_o} = -\frac{d_i}{d_o}\]
			
		
		
		% PUNTO 3
		\item \textbf{(25\%)}
		De acuerdo con las siguientes imágenes que representan el comportamiento de los rayos de luz al incidir sobre espejos esféricos, indique cuál es la opción \textbf{correcta} y justifique su elección:
		
		\vspace{0.3cm}
		\begin{center}
			\begin{tikzpicture}[scale=0.85, >=Stealth]
				% Opción a
				\begin{scope}[shift={(0,0)}]
					\draw[thick] (-1.5,-1.5) rectangle (1.5,1.5);
					\node at (-1.9,1.2) {\textbf{a.}};
					\draw[dashed] (-1.5,0) -- (1.5,0);
					
					% Espejo convexo
					\draw[thick] (0.4, 1.2) to[bend left=-55] (0.4, -1.2); 
					
					% Foco y Centro
					\node[below, font=\tiny] at (0.3, 0) {F};
					\draw (0.3, 0.05) -- (0.3, -0.05);
					\node[below, font=\tiny] at (1.0, 0) {C};
					\draw (1.0, 0.05) -- (1.0, -0.05);
					
					% Rayo Incidente (Intersección exacta)
					\draw[red] (-1.5, 0.6) -- (-0.08, 0.6);
					\draw[red, ->] (-1.5, 0.6) -- (-0.79, 0.6);
					
					% Rayo Virtual
					\draw[red, dashed] (-0.08, 0.6) -- (0.3, 0);
					
					% Rayo Reflejado
					\draw[red] (-0.08, 0.6) -- (-0.65, 1.5);
					\draw[red, ->] (-0.08, 0.6) -- (-0.365, 1.05);
				\end{scope}
				
				% Opción b
				\begin{scope}[shift={(4.5,0)}]
					\draw[thick] (-1.5,-1.5) rectangle (1.5,1.5);
					\node at (-1.9,1.2) {\textbf{b.}};
					\draw[dashed] (-1.5,0) -- (1.5,0);
					
					% Espejo convexo
					\draw[thick] (0.1, 1.2) to[bend left=-55] (0.1, -1.2); 
					
					% Foco y Centro
					\node[below, font=\tiny] at (0.3, 0) {F};
					\draw (0.3, 0.05) -- (0.3, -0.05);
					\node[below, font=\tiny] at (1.0, 0) {C};
					\draw (1.0, 0.05) -- (1.0, -0.05);
					
					% Rayo incidente 
					\draw[blue, ->] (-0.28, 0.8) -- (0.3, 0);
					\draw[blue, ->] (-1.5, 1.16) -- (-0.28, 0.8);
					
					% Rayo despues apuntando a F
					%\draw[blue] (-1.5, 1.16) -- (0.3, 0);
					%\draw[blue, ->] (-1.5, 1.16) -- (-0.6, 0.58);
				\end{scope}
				
				% Opción c
				\begin{scope}[shift={(9,0)}]
					\draw[thick] (-1.5,-1.5) rectangle (1.5,1.5);
					\node at (-1.9,1.2) {\textbf{c.}};
					\draw[dashed] (-1.5,0) -- (1.5,0);
					
					% Espejo cóncavo 
					\draw[thick] (0.14, 1.2) to[bend right=-55] (0.14, -1.2);
					
					% Foco y Centro 
					\node[below, font=\tiny] at (-0.5, 0) {C};
					\draw (-0.5, 0.05) -- (-0.5, -0.05);
					\node[below, font=\tiny] at (0.2, 0) {F};
					\draw (0.2, 0.05) -- (0.2, -0.05);
					
					% Rayo Incidente
					\draw[red] (-1.5, 0.6) -- (0.58, 0.6);
					\draw[red, ->] (-1.5, 0.6) -- (-0.46, 0.6);
					
					% Rayo Reflejado (erróneo a propósito)
					\draw[red] (0.58, 0.6) -- (1.5, 1.1);
					\draw[red, ->] (0.58, 0.6) -- (1.04, 0.85);
				\end{scope}
				
				% Opción d
				\begin{scope}[shift={(13.5,0)}]
					\draw[thick] (-1.5,-1.5) rectangle (1.5,1.5);
					\node at (-1.9,1.2) {\textbf{d.}};
					\draw[dashed] (-1.5,0) -- (1.5,0);
					
					% Espejo cóncavo
					\draw[thick] (0.14, 1.2) to[bend right=-55] (0.14, -1.2);
					
					% Foco y Centro
					\node[below, font=\tiny] at (-0.3, 0) {C};
					\draw (-0.3, 0.05) -- (-0.3, -0.05);
					\node[below, font=\tiny] at (0.4, 0) {F};
					\draw (0.4, 0.05) -- (0.4, -0.05);
					
					% Rayo incidente (pasa por C)
					\draw[blue] (-1.5, 0.63) -- (0.65, -0.5);
					\draw[blue, ->] (-1.5, 0.63) -- (-0.42, 0.06); 
					
					% Rayo reflejado (erróneo pasando por F)
					\draw[blue] (0.65, -0.5) -- (-0.35, 1.5);
					\draw[blue, ->] (0.65, -0.5) -- (0.15, 0.5);
				\end{scope}
			\end{tikzpicture}
		\end{center}
		
		
		
		% PUNTO 4
		\item \textbf{(25\%)}
		Responda Falso (\textbf{F}) o Verdadero (\textbf{V}) a cada una de las siguientes afirmaciones y justifique su respuesta con base en las leyes de la reflexión estudiadas en el TAU 5:
		
		\begin{enumerate}[label=\alph*.), itemsep=0.4cm]
			\item Los espejos planos producen únicamente imágenes virtuales, derechas y del mismo tamaño que el objeto, por lo que la afirmación ``\textit{Los espejos planos producen imágenes virtuales y reales}'' es: \underline{\hspace{1.5cm}} \\
			
			
			\item Al incidir sobre una superficie rugosa o irregular, los rayos de un haz de luz se reflejan en distintas direcciones, por lo que la afirmación ``\textit{Los rayos de un haz de luz no siempre se reflejan en una misma dirección al incidir sobre una superficie}'' es: \underline{\hspace{1.5cm}} \\
			
		\end{enumerate}
		
		
		
	\end{enumerate}
	
	\vfill
	\begin{center}
		\textit{\color{valdark!70} ``Hacia un Aprendizaje Significativo, Autónomo y en Libertad'' -- Liceo V.A.L.}
	\end{center}
	
\end{document}